\setcounter{section}{21}

  \zwischenueberschrift{Manifold berbatas}

\bild{ \begin{center}
\includegraphics[width=5.5cm]{ \bildeinlesungsvg {Runge_theorem} {svg} }
\end{center}
\bildtext {Suatu manifold berdimensi dua dengan batas. Batasnya terdiri atas empat busur tertutup.} }

\bildlizenz { Runge theorem.svg } {} {Oleg Alexandrov} {Commons} {PD} {}

\inputdefinition
{}
{

Misalkan

\mavergleichskette
{\vergleichskette
{ n
}
{ \in }{ \N
}
{  }{
}
{    }{
}
{   }{
}
}
{}{}{}
dan

\mavergleichskette
{\vergleichskette
{ k
}
{ \in }{ \overline{ \N }_+
}
{  }{
}
{    }{
}
{   }{
}
}
{}{}{.}
Suatu
\definitionsverweis {ruang topologis}{}{}
\definitionsverweis {Hausdorff}{}{}
$M$ bersama dengan suatu
\definitionsverweis {tutupan terbuka}{}{}

\mavergleichskette
{\vergleichskette
{ M
}
{ = }{ \bigcup_{i \in I} U_i
}
{  }{
}
{    }{
}
{   }{
}
}
{}{}{}
dan
\definitionsverweis {peta-peta}{}{}
\maabbdisp {\alpha_i} {U_i } { V_i
} {,}
dengan

\mavergleichskette
{\vergleichskette
{ V_i
}
{ \subseteq }{ H
}
{ \subset }{ \R^n
}
{    }{
}
{   }{
}
}
{}{}{}
merupakan himpunan-himpunan terbuka di dalam
\definitionsverweis {setengah ruang Euklides}{}{}
$H$ berdimensi $n$, serta dengan sifat bahwa
\definitionsverweis {pemetaan-pemetaan transisi}{}{}
\maabbdisp {\alpha_j \circ \alpha_i^{-1}} {V_i \cap \alpha_i(U_i \cap U_j) } { V_j \cap \alpha_j(U_i \cap U_j)
} {}
merupakan
$C^k$-\definitionsverweis {difeomorfisme}{}{}{,}
disebut \definitionswort {manifold-$C^k$ berbatas}{} atau \definitionswort {manifold diferensiabel berbatas}{}
\zusatzklammer {berderajat $k$} {} {,}
atau \definitionswort {manifold dengan batas}{.} Himpunan peta

\mathbed {(U_i,\alpha_i)} {}
 {i \in I} {}
 {} {} {} {,}
juga disebut \definitionswort {atlas-$C^k$}{} dari manifold berbatas tersebut.

}
Karena himpunan-himpunan terbuka di dalam setengah ruang $H$ yang sama sekali tidak memotong batas

\mavergleichskette
{\vergleichskette
{ \partial H
}
{ = }{ 0 \times \R^{n-1}
}
{  }{
}
{    }{
}
{   }{
}
}
{}{}{}
juga diperbolehkan, konsep ini mencakup konsep manifold diferensiabel. Batas suatu manifold berbatas
\zusatzklammer {yang segera akan kita definisikan dengan cara yang wajar} {} {}
dapat saja kosong. Hal ini terjadi tepat pada manifold diferensiabel \anfuehrung{biasa}{}.

\inputdefinition
{}
{

Misalkan $M$ suatu
\definitionsverweis {manifold berbatas}{}{.}
Maka \definitionswort {batas}{} dari $M$, yang ditulis
\mathl{\partial M}{,} didefinisikan oleh

\mavergleichskettedisp
{\vergleichskette
{ \partial M
}
{ =} { \left\{  x \in M  \mid   \alpha_i(x) \in  \partial H \cap V_i  \text{ untuk } \text{suatu } i \in I         \right\}
}
{ } {
}
{ } {
}
{ } {
}
}
{}{}{}
dengan $\alpha_i$ adalah peta-peta.

}
Pada definisi ini, peta mana pun dapat digunakan untuk menguji apakah suatu titik yang diberikan merupakan titik batas.
__NOEDITSECTION__

\inputfaktbeweis
{Mannigfaltigkeit mit Rand/Rand/Elementare Eigenschaften/Fakt}
{Lema}
{}
{

\faktsituation {Misalkan $M$ suatu
\definitionsverweis {manifold diferensiabel berbatas}{}{}

\mathl{\partial M}{.}}
\faktuebergang {Maka pernyataan-pernyataan berikut berlaku.}
\faktfolgerung {\aufzaehlungdrei{Suatu titik

\mavergleichskette
{\vergleichskette
{ P
}
{ \in }{ M
}
{  }{
}
{    }{
}
{   }{
}
}
{}{}{}
merupakan titik batas jika dan hanya jika hal ini berlaku bagi setiap
\definitionsverweis {peta}{}{,}
yang domain petanya memuat titik tersebut\zusatzfussnote {Pada manifold topologis berbatas, yaitu ketika pemetaan-pemetaan transisinya hanya berupa homeomorfisme, pernyataan ini juga berlaku, tetapi dengan pembuktian yang berbeda.} {} {.}
}{Batas $\partial M$ merupakan himpunan
\definitionsverweis {tertutup}{}{.}
}{Komplemen
\mathl{M \setminus \partial M}{} merupakan suatu
\definitionsverweis {manifold diferensiabel}{}{}
\zusatzklammer {tanpa batas} {} {.}
}}
\faktzusatz {}
\faktzusatz {}

}
{

\teilbeweis {}{}{}
{(1). Misalkan

\mavergleichskette
{\vergleichskette
{ P
}
{ \in }{ M
}
{  }{
}
{    }{
}
{   }{
}
}
{}{}{}
dan

\mavergleichskette
{\vergleichskette
{ P
}
{ \in }{ U
}
{  }{
}
{    }{
}
{   }{
}
}
{}{}{}
suatu domain peta dengan dua peta
\maabbdisp {\alpha_1} { U } { V_1
} {}
dan
\maabbdisp {\alpha_2} { U } { V_2
} {}
dengan
\mathkor {} {V_1 \subseteq H_1} {dan} {V_2 \subseteq H_2} {}
terbuka di dalam
\definitionsverweis {setengah ruang Euklides}{}{}

\mavergleichskette
{\vergleichskette
{ H_1, H_2
}
{ \cong }{ \R_{\geq 0} \times \R^{n-1}
}
{  }{
}
{    }{
}
{   }{
}
}
{}{}{.}
Pemetaan pergantian peta

\mavergleichskette
{\vergleichskette
{ \varphi
}
{ = }{ \alpha_2 \circ \alpha_1^{-1}
}
{  }{
}
{    }{
}
{   }{
}
}
{}{}{}
merupakan suatu
\definitionsverweis {difeomorfisme}{}{,}
dan menurut [[Halbraum/Offene Menge/Diffeomorphismus/Offene Umgebung/Aufgabe|Soal 21.14]] ini berarti bahwa bagi setiap titik

\mavergleichskette
{\vergleichskette
{ Q
}
{ \in }{ V_1
}
{  }{
}
{    }{
}
{   }{
}
}
{}{}{,}
terdapat lingkungan-lingkungan terbuka
\mathkor {} {W_1} {dan} {W_2} {}
di $\R^n$ dengan

\mavergleichskette
{\vergleichskette
{ Q
}
{ \in }{ W_1
}
{  }{
}
{    }{
}
{   }{
}
}
{}{}{}
serta suatu perluasan
\definitionsverweis {difeomorfik}{}{}
\maabbdisp {\tilde{\varphi}} {W_1 } { W_2
} {}
dari
\mathl{\varphi\!\mid_{W_1 \cap H_1 }}{.} Oleh karena itu,
\mathl{\tilde{\varphi} \left(  H_1^+ \cap W_1  \right)}{} terbuka di $W_2$.

Sekarang, tetapkan

\mavergleichskette
{\vergleichskette
{ Q_1
}
{ = }{ \alpha_1(P)
}
{  }{
}
{    }{
}
{   }{
}
}
{}{}{}
dan pilih
\mathl{W_1,W_2, \tilde{\varphi}}{} dengan sifat-sifat yang baru saja disebutkan. Jika $Q_1$ bukan titik batas dalam peta pertama, maka

\mavergleichskette
{\vergleichskette
{ Q_1
}
{ \in }{ H_1^+ \cap W_1
}
{  }{
}
{    }{
}
{   }{
}
}
{}{}{}
merupakan suatu lingkungan terbuka dan karenanya

\mavergleichskette
{\vergleichskette
{ \alpha_2 (P)
}
{ \in }{ \tilde{\varphi} \left(  H_1^+ \cap W_1  \right)
}
{  }{
}
{    }{
}
{   }{
}
}
{}{}{}
merupakan suatu lingkungan terbuka di

\mavergleichskette
{\vergleichskette
{ W_2
}
{ \subseteq }{ \R^n
}
{  }{
}
{    }{
}
{   }{
}
}
{}{}{.}
Selain itu,

\mavergleichskette
{\vergleichskette
{ \tilde{\varphi} \left(  H_1^+ \cap W_1  \right)
}
{ \subseteq }{ H_2
}
{  }{
}
{    }{
}
{   }{
}
}
{}{}{.}
Dengan kata lain,

\mavergleichskette
{\vergleichskette
{ \alpha_2(P)
}
{ \in }{ H_2
}
{  }{
}
{    }{
}
{   }{
}
}
{}{}{}
memiliki suatu lingkungan yang terbuka di $\R^n$ dan termuat di dalam $H_2$; maka menurut [[Halbraum/Randpunkt/Keine offene Menge im Raum/Aufgabe|Soal 21.13]], titik ini juga tidak mungkin menjadi titik batas dalam peta kedua.}
{}
\teilbeweis {}{}{}
{(2). Misalkan

\mavergleichskette
{\vergleichskette
{ P
}
{ \notin }{ \partial M
}
{  }{
}
{    }{
}
{   }{
}
}
{}{}{}
dan misalkan

\mavergleichskette
{\vergleichskette
{ P
}
{ \in }{ U
}
{  }{
}
{    }{
}
{   }{
}
}
{}{}{}
suatu domain peta dengan
\definitionsverweis {homeomorfisme}{}{}
\maabbdisp {\alpha} { U } { V
} {}
dengan

\mavergleichskette
{\vergleichskette
{ V
}
{ \subseteq }{ \R_{\geq 0} \times \R^{n-1}
}
{ = }{ H
}
{    }{
}
{   }{
}
}
{}{}{}
terbuka. Karena $P$ bukan titik batas, komponen pertama
\mathl{\alpha(P)}{} positif, sehingga terdapat suatu himpunan terbuka

\mavergleichskette
{\vergleichskette
{ \alpha(P)
}
{ \in }{ V'
}
{ \subseteq }{ \R_{+} \times \R^{n-1}
}
{ =   }{ H_+
}
{   }{
}
}
{}{}{.}
Dengan demikian,

\mavergleichskette
{\vergleichskette
{ U'
}
{ = }{ \alpha^{-1}(V')
}
{  }{
}
{    }{
}
{   }{
}
}
{}{}{}
merupakan suatu lingkungan terbuka dari $P$ yang
\zusatzklammer {menurut Bagian 1} {} {}
tidak memotong batas.}
{}
\teilbeweis {}{}{}
{(3). Bagi setiap titik

\mavergleichskette
{\vergleichskette
{ P
}
{ \notin }{ \partial M
}
{  }{
}
{    }{
}
{   }{
}
}
{}{}{}
kita dapat menentukan, seperti pada (2), suatu domain peta yang terpisah dari batas dan citra petanya merupakan himpunan terbuka di $\R^n$. Jadi, kita memperoleh suatu
\definitionsverweis {manifold}{}{}
\zusatzklammer {tanpa batas} {} {.}}
{}

}

Konsep pemetaan diferensiabel, difeomorfisme, dan ruang singgung juga terbawa ke manifold berbatas; lihat
[[Mannigfaltigkeit mit Rand/Differenzierbare Abbildung/Tangentialraum/Begriffseinführung/Aufgabe|Soal 21.2]].

__NOEDITSECTION__

\inputfaktbeweis
{Mannigfaltigkeit mit Rand/Rand ist Mannigfaltigkeit/Fakt}
{Teorema}
{}
{

\faktsituation {Misalkan $M$ suatu
\definitionsverweis {manifold diferensiabel berbatas}{}{} berdimensi $n$.}
\faktfolgerung {Maka batas
\mathl{\partial M}{} merupakan suatu
\definitionsverweis {manifold diferensiabel}{}{}
\zusatzklammer {tanpa batas} {} {} berdimensi
\mathl{n-1}{.}}
\faktzusatz {}
\faktzusatz {}

}
{

Pertama-tama,

\mavergleichskette
{\vergleichskette
{ L
}
{ = }{ \partial M
}
{  }{
}
{    }{
}
{   }{
}
}
{}{}{}
dengan
\definitionsverweis {topologi terinduksi}{}{}
merupakan suatu
\definitionsverweis {ruang Hausdorff}{}{.}
Misalkan

\mavergleichskette
{\vergleichskette
{ P
}
{ \in }{ L
}
{  }{
}
{    }{
}
{   }{
}
}
{}{}{}
dan misalkan
\maabbdisp {\alpha} { U } { V
} {}
suatu
\definitionsverweis {peta}{}{}
dengan

\mavergleichskette
{\vergleichskette
{ V
}
{ \subseteq }{ H
}
{  }{
}
{    }{
}
{   }{
}
}
{}{}{}
terbuka dan

\mavergleichskette
{\vergleichskette
{ \alpha(P)
}
{ \in }{ \partial H
}
{  }{
}
{    }{
}
{   }{
}
}
{}{}{.}
Karena $\alpha$ merupakan suatu
\definitionsverweis {homeomorfisme}{}{}
dan menurut
[[Mannigfaltigkeit mit Rand/Rand/Elementare Eigenschaften/Fakt|Lema 21.3 (1)]]
titik batas berkorespondensi dengan titik batas pada setiap peta, peta ini menginduksi suatu homeomorfisme
\maabbdisp {} {U \cap L} {V \cap \partial H
} {.}
Di sini,
\mathl{U \cap L}{} merupakan suatu lingkungan terbuka dari $P$ di $L$, sehingga himpunan-himpunan ini dapat kita gunakan sebagai domain peta. Sekarang tinjau suatu pergantian peta. Kita boleh langsung memulai dengan satu domain peta

\mavergleichskette
{\vergleichskette
{ U
}
{ \subseteq }{ M
}
{  }{
}
{    }{
}
{   }{
}
}
{}{}{}
dan dua peta
\maabb {\alpha_1} { U } { V_1
} {} dan \maabb {\alpha_2} { U } { V_2
} {}
dengan

\mavergleichskette
{\vergleichskette
{ V_1,V_2
}
{ \subseteq }{ H
}
{  }{
}
{    }{
}
{   }{
}
}
{}{}{}
terbuka. Dengan demikian, terdapat suatu
$C^1$-\definitionsverweis {difeomorfisme
}{}{}
\maabbdisp {\varphi=\alpha_2 \circ \alpha_1^{-1}} {V_1 } { V_2
} {.}
Pertama, hal ini berarti terdapat suatu homeomorfisme
\maabb {} { \partial H \cap V_1 } { \partial H \cap V_2
} {.}
Sifat difeomorfisme dari $\varphi$ berarti bahwa bagi setiap titik

\mavergleichskette
{\vergleichskette
{ P
}
{ \in }{ U \cap L
}
{  }{
}
{    }{
}
{   }{
}
}
{}{}{,}
terdapat lingkungan-lingkungan terbuka

\mavergleichskette
{\vergleichskette
{ \alpha_1(P)
}
{ \in }{ W_1
}
{ \subseteq }{ \R^n
}
{    }{
}
{   }{
}
}
{}{}{}
dan

\mavergleichskette
{\vergleichskette
{ \alpha_2(P)
}
{ \in }{ W_2
}
{ \subseteq }{ \R^n
}
{    }{
}
{   }{
}
}
{}{}{}
serta suatu
\definitionsverweis {perluasan}{}{}
\maabbdisp {\tilde{\varphi}} { W_1 } { W_2
} {}
yang difeomorfik dari $\varphi$, dari
\mathkor {} {V_1 \cap W_1} {ke} {V_2 \cap W_2} {.}
Menurut [[Diffeomorphismus/Halbraum auf Halbraum/Induziert Diffeomorphismus/Aufgabe|Soal 21.19]], perluasan ini juga menginduksi suatu
$C^1$-\definitionsverweis {difeomorfisme}{}{}
antara batas-batas
\mathl{\partial H \cap W_1}{} dan
\mathl{\partial H \cap W_2}{,} sehingga secara keseluruhan terdapat suatu difeomorfisme
\maabbdisp {\varphi \vert_{\partial H}} { \partial H \cap V_1 } { \partial H \cap V_2
} {.}

}

Kita sudah mengetahui bahwa serat suatu pemetaan reguler diferensiabel antara manifold-manifold diferensiabel memiliki struktur submanifold tertutup. Teorema berikut didasarkan pada argumen serupa dan menjamin keberadaan sangat banyak manifold berbatas.
__NOEDITSECTION__

\inputfaktbeweis
{Reguläre Funktion auf Mannigfaltigkeit/Urbild halbseitiger Intervalle/Mannigfaltigkeit mit Rand/Fakt}
{Teorema}
{}
{

\faktsituation {Misalkan $L$ suatu
\definitionsverweis {manifold diferensiabel}{}{}
dan misalkan
\maabbdisp {f} { L } {\R
} {}
suatu
\definitionsverweis {fungsi yang terdiferensialkan secara kontinu}{}{.}
Misalkan

\mavergleichskette
{\vergleichskette
{ a
}
{ \in }{ \R
}
{  }{
}
{    }{
}
{   }{
}
}
{}{}{.}}
\faktvoraussetzung {Setiap titik dalam
\definitionsverweis {serat}{}{}

\mavergleichskette
{\vergleichskette
{ F
}
{ = }{ f^{-1}(a)
}
{  }{
}
{    }{
}
{   }{
}
}
{}{}{}
di atas $a$ merupakan titik
\definitionsverweis {reguler}{}{.}}
\faktfolgerung {Maka himpunan-himpunan bagian

\mathdisp {M= \left\{  x \in L  \mid   f(x) \leq a        \right\} \text{   dan } N= \left\{  x \in L  \mid   f(x) \geq a        \right\}} {  }

merupakan
\definitionsverweis {manifold diferensiabel berbatas}{}{,}
dan \definitionsverweis {batas}{}{} masing-masing sama dengan $F$.}
\faktzusatz {}
\faktzusatz {}

}
{

Untuk menyederhanakan notasi, misalkan

\mavergleichskette
{\vergleichskette
{ a
}
{ = }{ 0
}
{  }{
}
{    }{
}
{   }{
}
}
{}{}{,}
dan kita batasi perhatian pada $M$. Kita mempunyai

\mavergleichskettedisp
{\vergleichskette
{ M
}
{ =} { \left\{  x\in L  \mid  f(x)<0        \right\}  \uplus F
}
{ } {
}
{ } {
}
{ } {
}
}
{}{}{,}
dengan himpunan di sebelah kiri merupakan suatu
\definitionsverweis {himpunan terbuka}{}{}
di $L$ dan dengan demikian merupakan suatu submanifold terbuka. Jadi, hal yang menentukan ialah menunjukkan bahwa bagi titik-titik dari serat $F$ terdapat peta-peta, dan dengan demikian suatu struktur manifold. Misalkan

\mavergleichskette
{\vergleichskette
{ P
}
{ \in }{ F
}
{  }{
}
{    }{
}
{   }{
}
}
{}{}{.}
Menurut pembuktian
[[Satz über implizite Abbildungen/R/Fakt|teorema pemetaan implisit]],
bagi

\mavergleichskette
{\vergleichskette
{ P
}
{ \in }{ F
}
{  }{
}
{    }{
}
{   }{
}
}
{}{}{}
terdapat suatu lingkungan
\zusatzklammer {koordinat} {} {}
terbuka

\mavergleichskette
{\vergleichskette
{ P
}
{ \in }{ U
}
{ \subseteq }{ L
}
{    }{
}
{   }{
}
}
{}{}{}
dan suatu peta
\maabbdisp {\alpha} { U } { V
} {,}

\mavergleichskette
{\vergleichskette
{ V
}
{ \subseteq }{ \R^n
}
{  }{
}
{    }{
}
{   }{
}
}
{}{}{,}
sedemikian sehingga

\mavergleichskette
{\vergleichskette
{ f
}
{ = }{ x_1 \circ \alpha
}
{  }{
}
{    }{
}
{   }{
}
}
{}{}{.}
Dalam hal ini,
\mathl{F \cap U}{} berkorespondensi dengan
\mathl{\left\{  x \in V  \mid   x_1 = 0         \right\}}{} dan
\mathl{M \cap U}{} berkorespondensi dengan
\mathl{H_{\leq 0} \cap V}{,} sehingga pembatasan $\alpha$ pada $M$ menghasilkan suatu peta untuk $M$ di $P$. Pergantian-pergantian petanya
$C^1$-\definitionsverweis {difeomorfik}{}{,}
karena hal ini berlaku bagi peta-peta
\zusatzklammer {penuh} {} {}
pada $L$.

}

\inputbeispiel{}
{

Untuk $r>0$, \definitionsverweis {bola tertutup}{}{}

\mavergleichskettedisp
{\vergleichskette
{ B  \left(  0 ,r \right)
}
{ =} { \left\{  x \in \R^n   \mid   \sqrt{x_1^2 + \cdots + x_n^2} \leq r         \right\}
}
{ } {
}
{ } {
}
{ } {
}
}
{}{}{}
merupakan suatu
$C^\infty$-\definitionsverweis {manifold diferensiabel}{}{}
dengan \stichwort {sfer} {}

\mavergleichskettedisp
{\vergleichskette
{ S \left(   0 , r  \right)
}
{ =} { \left\{  x \in \R^n   \mid   \sqrt{x_1^2 + \cdots + x_n^2} =  r         \right\}
}
{ } {
}
{ } {
}
{ } {
}
}
{}{}{}
sebagai
\definitionsverweis {batas}{}{.}
Hal ini langsung mengikuti
[[Reguläre Funktion auf Mannigfaltigkeit/Urbild halbseitiger Intervalle/Mannigfaltigkeit mit Rand/Fakt|Teorema 21.5]]
yang diterapkan pada fungsi mulus
\maabbeledisp {g} {\R^n } {\R
} {(x_1 , \ldots , x_n) } { x_1^2 + \cdots + x_n^2
} {,}
dengan nilai reguler $r^2$; setiap titik $x\in g^{-1}(r^2)$ memenuhi
$x\neq 0$, sehingga $g$ bersifat
\definitionsverweis {reguler}{}{} pada seluruh serat itu.

}

\inputbeispiel{}
{

Suatu
\definitionsverweis {balok}{}{} \definitionsverweis {tertutup}{}{}

\mathdisp {[a_1,b_1] \times [a_2,b_2] \times \cdots \times [a_n,b_n] \subseteq \R^n} {  }

untuk
\mathl{n \geq 2}{} bukan merupakan suatu
\definitionsverweis {manifold diferensiabel berbatas}{}{,}
karena balok itu tidak hanya mempunyai sisi, tetapi juga
\zusatzklammer {bergantung pada dimensinya} {} {}
sudut dan rusuk. Suatu persegi panjang mempunyai empat titik sudut yang tidak dapat diberi struktur manifold diferensiabel berbatas
\zusatzklammer {yang kompatibel dengan ruang sekitar} {} {}
\zusatzklammer {karena persegi panjang tertutup
\definitionsverweis {homeomorfik}{}{}
dengan cakram tertutup, suatu struktur manifold diferensiabel berbatas dapat didefinisikan padanya; tetapi dengan struktur ini, penyematan alami persegi panjang tersebut ke dalam $\R^2$ tidak diferensiabel} {} {,}
sedangkan suatu balok berdimensi tiga mempunyai dua belas rusuk dan delapan sudut yang tidak memiliki struktur manifold alami.

Namun, jika sudut, rusuk, dan seterusnya tersebut dihapus dan hanya \anfuehrung{sisi-sisi berkodimensi satu}{} dipertahankan, kita memperoleh suatu
\zusatzklammer {tak kompak} {} {}
manifold berbatas. Batasnya merupakan gabungan saling lepas dari muka-muka hiper tersebut. Seperti pada setiap manifold berbatas, batas ini tertutup di dalam manifold, tetapi tidak tertutup di dalam ruang Euklides sekitar.

}

  \zwischenueberschrift{Orientasi pada manifold berbatas}

Misalkan pada $\R^n$ vektor-vektor standar
\mathl{e_1 , \ldots , e_n}{} memberikan suatu
\definitionsverweis {orientasi}{}{;}
selanjutnya pilih setengah ruang

\mavergleichskettedisp
{\vergleichskette
{ H_{\leq 0}
}
{ =} { \left\{  x \in \R^n   \mid   x_1 \leq 0         \right\}
}
{ } {
}
{ } {
}
{ } {
}
}
{}{}{}
sebagai \anfuehrung{setengah ruang dalam}{}. Orientasi pada hiperbidang
\zusatzklammer {yaitu batas manifold berbatas \mathlk{H_{\leq 0}}{}} {} {}

\mavergleichskettedisp
{\vergleichskette
{ E
}
{ =} { \left\{  x \in \R^n   \mid   x_1 = 0         \right\}
}
{ } {
}
{ } {
}
{ } {
}
}
{}{}{}
yang didefinisikan oleh basis
\mathl{e_2 , \ldots , e_n}{} disebut \stichwort {orientasi oleh normal luar} {.} Suatu basis sembarang
\mathl{v_2 , \ldots , v_n}{} dari $E$ merepresentasikan orientasi ini jika dan hanya jika, bagi suatu vektor sembarang

\mavergleichskette
{\vergleichskette
{ v
}
{ \in }{ H_+
}
{  }{
}
{    }{
}
{   }{
}
}
{}{}{}
\zusatzklammer {artinya, mengarah ke \anfuehrung{luar}{,} yakni keluar dari setengah ruang} {} {}
basis
\mathl{v,v_2 , \ldots , v_n}{}
\zusatzklammer {jadi $v$ ditempatkan pertama} {} {}
dari $\R^n$ merepresentasikan orientasi semula
\zusatzfussnote {Untuk suatu setengah garis

\mavergleichskette
{\vergleichskette
{H
}
{ = }{ \R_{\geq 0}
}
{ \subseteq }{ \R
}
{    }{
}
{   }{
}
}
{}{}{}
dengan satu-satunya titik batas $\{0\}$, hal ini harus ditafsirkan sebagai berikut. Kedua orientasi pada
\mathl{\{0\}}{} adalah
\mathkor {} {+} {dan} {-} {,}
dan $-$ merepresentasikan orientasi oleh normal luar, sebab bagi suatu vektor yang mengarah ke luar

\mavergleichskette
{\vergleichskette
{ w
}
{ \in }{ \R_-
}
{  }{
}
{    }{
}
{   }{
}
}
{}{}{}
vektor lawannya $-w$ merepresentasikan orientasi standar dari $\R$. Sebaliknya, untuk setengah ruang negatif
\mathl{\R_{\leq 0}}{} di titik nol, $+$ merepresentasikan orientasi oleh normal luar} {.} {.}

Hubungan antara orientasi pada ruang vektor real dan orientasi pada batas suatu setengah ruang ini terbawa ke manifold berbatas. Hal yang penting di sini ialah bahwa ruang singgung
\mathl{T_PM}{} di suatu titik batas $P$ memuat suatu hiperbidang kanonik, yaitu ruang singgung
\mathl{T_P (\partial M)}{} dari batas. Dalam hal ini manifold tersebut menentukan suatu \anfuehrung{separuh dalam}{} dan suatu \anfuehrung{separuh luar}{} dari ruang singgung.

__NOEDITSECTION__

\inputfaktbeweis
{Mannigfaltigkeit mit Rand/Orientierung/Randorientierung/Fakt}
{Teorema}
{}
{

\faktsituation {Misalkan $M$ suatu
\definitionsverweis {manifold diferensiabel berbatas}{}{,}}
\faktvoraussetzung {yang dilengkapi dengan suatu
\definitionsverweis {orientasi}{}{.}}
\faktfolgerung {Maka
\definitionsverweis {manifold batas}{}{} juga dilengkapi dengan suatu orientasi kanonik, yaitu orientasi yang pada setiap
\definitionsverweis {peta}{}{} ditentukan oleh
\definitionsverweis {normal luar}{}{.}}
\faktzusatz {}
\faktzusatz {}

}
{

Untuk setiap peta
\zusatzklammer {berorientasi} {} {}
\maabbdisp {\alpha} { U } { V
} {}
dengan

\mavergleichskette
{\vergleichskette
{ U
}
{ \subseteq }{ M
}
{  }{
}
{    }{
}
{   }{
}
}
{}{}{}
terbuka, peta terinduksi
\maabbdisp {} {U \cap \partial M } { V \cap \partial H
} {}
dilengkapi dengan
\definitionsverweis {orientasi oleh normal luar}{}{}
pada
\mathl{\partial H}{.} Menurut asumsi, semua pergantian peta pada $M$ mempunyai
\definitionsverweis {determinan fundamental}{}{}
positif di setiap titik terhadap basis-basis yang merepresentasikan orientasi; kita harus menunjukkan bahwa hal ini juga berlaku bagi pergantian peta yang terinduksi. Kita dapat memulai dengan suatu domain peta terbuka

\mavergleichskette
{\vergleichskette
{ P
}
{ \in }{ U
}
{ \subseteq }{ M
}
{    }{
}
{   }{
}
}
{}{}{}
dan dua peta
\maabbdisp {\alpha_1} { U } { V_1
} {}
dan
\maabbdisp {\alpha_2} { U } { V_2
} {}
serta meninjau pemetaan transisi
\maabbdisp {\varphi = \alpha_2 \circ \alpha_1^{-1}} { V_1 } { V_2
} {}
dengan himpunan-himpunan terbuka

\mavergleichskette
{\vergleichskette
{ V_1
}
{ \subseteq }{ H_1
}
{ \subset }{ \R^n
}
{    }{
}
{   }{
}
}
{}{}{}
dan

\mavergleichskette
{\vergleichskette
{ V_2
}
{ \subseteq }{ H_2
}
{ \subset }{ \R^n
}
{    }{
}
{   }{
}
}
{}{}{.}
Misalkan
\mathl{v_2 , \ldots , v_n}{} suatu basis dari
\mathl{\partial H_1}{} dan

\mavergleichskette
{\vergleichskette
{ v_1
}
{ \in }{ \R^n
}
{  }{
}
{    }{
}
{   }{
}
}
{}{}{}
sedemikian sehingga $v_1$ merepresentasikan normal luar dari $H_1$, yakni
\mathl{v_1, v_2 , \ldots , v_n}{} merepresentasikan orientasi $\R^n$
\zusatzklammer {misalkan $x_1 , \ldots , x_n$ adalah koordinat-koordinat yang bersesuaian} {} {;}
demikian pula,

\mavergleichskette
{\vergleichskette
{ w_1, w_2 , \ldots , w_n
}
{ \in }{ \R^n
}
{  }{
}
{    }{
}
{   }{
}
}
{}{}{}
harus memenuhi sifat-sifat yang bersesuaian. Kita tuliskan matriks fundamental dari $\varphi$ terhadap basis-basis ini, yakni matriks dengan entri-entri

\mathdisp {\left(  { \frac{   \partial  \varphi_i   }{  \partial   x_j   } } ( Q)  \right)_{1 \leq i,j \leq n}} { . }

Menurut asumsi, determinannya positif. Karena

\mavergleichskette
{\vergleichskette
{ \varphi ( \partial H_1)
}
{ \subseteq }{ \partial H_2
}
{  }{
}
{    }{
}
{   }{
}
}
{}{}{}
maka untuk suatu titik

\mavergleichskette
{\vergleichskette
{ Q
}
{ \in }{ \partial H_1
}
{  }{
}
{    }{
}
{   }{
}
}
{}{}{}
berlaku hubungan

\mavergleichskettedisp
{\vergleichskette
{ { \frac{   \partial  \varphi_1   }{  \partial   x_j   } } ( Q)
}
{ =} { 0
}
{ } {
}
{ } {
}
{ } {
}
}
{}{}{}
untuk

\mavergleichskette
{\vergleichskette
{ j
}
{ = }{ 2 , \ldots , n
}
{  }{
}
{    }{
}
{   }{
}
}
{}{}{.}
Menurut
[[Determinante/Körper/Entwicklungssatz/Fakt|teorema ekspansi]],
oleh karena itu tanda determinan matriks

\mathdisp {\left(  { \frac{   \partial  \varphi_i   }{  \partial   x_j   } } ( Q)  \right)_{ 2 \leq i,j \leq n}} { , }

yang merupakan matriks fundamental dari pemetaan transisi peta-peta batas
\maabbdisp {} {V_1 \cap \partial H_1 } { V_2 \cap \partial H_2
} {}
\zusatzklammer {terhadap basis
\mathl{v_2 , \ldots , v_n}{} dan
\mathl{w_2 , \ldots , w_n}{}} {} {}
di titik $Q$, hanya bergantung pada
\mathl{{ \frac{   \partial  \varphi_1   }{  \partial   x_1   } } ( Q)}{.} Dengan

\mavergleichskette
{\vergleichskette
{ Q
}
{ = }{ (0,a_2 , \ldots , a_n)
}
{  }{
}
{    }{
}
{   }{
}
}
{}{}{}
menurut
[[Differenzierbarkeit/K/Zusammenhang zwischen partieller Ableitung und Richtungsableitung/Fakt|Lema 44.3 (Analisis (Osnabrück 2021--2023))]],
kita memperoleh hubungan

\mavergleichskettedisp
{\vergleichskette
{ { \frac{   \partial  \varphi_1   }{  \partial   x_1   } } ( Q)
}
{ =} { \operatorname{lim}_{   \epsilon \rightarrow  0   } \,  { \frac{   \varphi_1 (\epsilon, a_2 , \ldots , a_n)  }{  \epsilon } }
}
{ } {
}
{ } {
}
{ } {
}
}
{}{}{.}
Karena $v_1$ merepresentasikan normal luar, untuk nilai negatif
\zusatzklammer {dengan nilai mutlak yang cukup kecil} {} {}
$\epsilon$, vektor dengan koordinat

\mavergleichskette
{\vergleichskette
{ (\epsilon,a_2 , \ldots , a_n)
}
{ \in }{ H_1
}
{  }{
}
{    }{
}
{   }{
}
}
{}{}{.}
Karena itu, vektor citranya harus termasuk dalam $H_2$, dan dengan demikian

\mavergleichskette
{\vergleichskette
{ \varphi_1(\epsilon,a_2 , \ldots , a_n)
}
{ < }{ 0
}
{  }{
}
{    }{
}
{   }{
}
}
{}{}{.}
Jadi, hasil bagi ini
\mathl{\geq 0}{,} dan hal yang sama berlaku bagi limitnya. Karena terdapat suatu difeomorfisme, limit tersebut bahkan harus positif, dan pernyataan pun terbukti.

}

 [[Kategorie:Latexseite]]
